\subsection{Scaled Multiplicative Gaussian Model}

We consider a Gaussian observation model where the variance scales linearly with the mean, similar to the Poisson model but maintaining the Gaussian framework:
\begin{equation}
Y_i(t) = \mu_i(t) + \varepsilon_i(t)
\end{equation}
where $\mu_i(t) = \beta(t) S_i(t) \sum_j C_{ij} I_j(t)$ is the mean incidence for region $i$ at time $t$, and $\varepsilon_i(t) \sim \mathcal{N}(0, \mu_i(t))$ are independent noise terms with variance proportional to the mean.

The observation vector is $\mathbf{Y} = [Y_1(0), Y_2(0), \ldots, Y_1(T-1), Y_2(T-1)]^T \in \mathbb{R}^{2T}$, with mean vector:
\begin{equation}
\boldsymbol{\mu}(\boldsymbol{\phi}) = [\mu_1(0), \mu_2(0), \ldots, \mu_1(T-1), \mu_2(T-1)]^T
\end{equation}

\subsubsection{Statistical Properties}

Under this scaled multiplicative model:
\begin{align}
\mathbb{E}[Y_i(t)] &= \mu_i(t) \\
\text{Var}[Y_i(t)] &= \mu_i(t)
\end{align}

This gives the same mean-variance relationship as the Poisson model, but maintains the Gaussian distributional assumption.

\subsubsection{Likelihood Function}

The observations are independently distributed:
\begin{equation}
Y_i(t) \sim \mathcal{N}(\mu_i(t), \mu_i(t))
\end{equation}

The likelihood function is:
\begin{equation}
f(\mathbf{Y}|\boldsymbol{\phi}) = \prod_{t=0}^{T-1} \prod_{i=1}^{2} \frac{1}{\sqrt{2\pi\mu_i(t)}} \exp\left(-\frac{(Y_i(t) - \mu_i(t))^2}{2\mu_i(t)}\right)
\end{equation}

\subsubsection{Log-Likelihood Function}

Taking the logarithm:
\begin{align}
\ell(\boldsymbol{\phi}) &= \log f(\mathbf{Y}|\boldsymbol{\phi}) \\
&= \sum_{t=0}^{T-1} \sum_{i=1}^{2} \left[-\frac{1}{2}\log(2\pi\mu_i(t)) - \frac{(Y_i(t) - \mu_i(t))^2}{2\mu_i(t)}\right] \\
&= -T\log(2\pi) - \frac{1}{2}\sum_{t=0}^{T-1}\sum_{i=1}^{2} \log\mu_i(t) - \frac{1}{2}\sum_{t=0}^{T-1}\sum_{i=1}^{2} \frac{(Y_i(t) - \mu_i(t))^2}{\mu_i(t)}
\end{align}

\subsubsection{Score Functions}

Taking the partial derivative with respect to parameter $\phi_k \in \{\theta, S_1(0), I_1(0), S_2(0), I_2(0)\}$:
\begin{align}
\frac{\partial \ell}{\partial \phi_k} &= -\frac{1}{2}\sum_{t=0}^{T-1}\sum_{i=1}^{2} \frac{1}{\mu_i(t)} \frac{\partial \mu_i(t)}{\partial \phi_k} + \frac{1}{2}\sum_{t=0}^{T-1}\sum_{i=1}^{2} \frac{\partial}{\partial \phi_k}\left[\frac{(Y_i(t) - \mu_i(t))^2}{\mu_i(t)}\right] \\
&= -\frac{1}{2}\sum_{t=0}^{T-1}\sum_{i=1}^{2} \frac{1}{\mu_i(t)} \frac{\partial \mu_i(t)}{\partial \phi_k} + \frac{1}{2}\sum_{t=0}^{T-1}\sum_{i=1}^{2} \left[\frac{-2(Y_i(t) - \mu_i(t))}{\mu_i(t)} + \frac{(Y_i(t) - \mu_i(t))^2}{\mu_i(t)^2}\right] \frac{\partial \mu_i(t)}{\partial \phi_k} \\
&= \sum_{t=0}^{T-1}\sum_{i=1}^{2} \frac{\partial \mu_i(t)}{\partial \phi_k} \underbrace{\left[-\frac{1}{2\mu_i(t)} - \frac{(Y_i(t) - \mu_i(t))}{\mu_i(t)} + \frac{(Y_i(t) - \mu_i(t))^2}{2\mu_i(t)^2}\right]}_{A_i(t)}
\end{align}

\subsubsection{Fisher Information Matrix}

To compute the Fisher Information Matrix using the outer product of gradients method, we need to evaluate:
\begin{equation}
[\mathbf{J}]_{kl} = \mathbb{E}\left[\frac{\partial \ell}{\partial \phi_k}\frac{\partial \ell}{\partial \phi_l}\right] = \sum_{t,i} \sum_{s,j} \mathbb{E}[A_i(t) A_j(s)] \frac{\partial \mu_i(t)}{\partial \phi_k}\frac{\partial \mu_j(s)}{\partial \phi_l}
\end{equation}

Since the noise terms are independent, $\mathbb{E}[A_i(t) A_j(s)] = 0$ for $(i,t) \neq (j,s)$. For $(i,t) = (j,s)$, we need to compute $\mathbb{E}[A_i(t)^2]$.

Let $Z_i(t) = Y_i(t) - \mu_i(t) \sim \mathcal{N}(0, \mu_i(t))$. Then:
\begin{equation}
A_i(t) = -\frac{1}{2\mu_i(t)} - \frac{Z_i(t)}{\mu_i(t)} + \frac{Z_i(t)^2}{2\mu_i(t)^2}
\end{equation}

Expanding $A_i(t)^2$:
\begin{align}
A_i(t)^2 &= \frac{1}{4\mu_i(t)^2} + \frac{Z_i(t)}{\mu_i(t)^2} + \frac{Z_i(t)^2}{4\mu_i(t)^3} + \frac{Z_i(t)^2}{\mu_i(t)^2} - \frac{Z_i(t)^3}{\mu_i(t)^3} + \frac{Z_i(t)^4}{4\mu_i(t)^4}
\end{align}

Using the moments of $Z_i(t) \sim \mathcal{N}(0, \mu_i(t))$:
\begin{itemize}
\item $\mathbb{E}[Z_i(t)] = 0$
\item $\mathbb{E}[Z_i(t)^2] = \mu_i(t)$
\item $\mathbb{E}[Z_i(t)^3] = 0$
\item $\mathbb{E}[Z_i(t)^4] = 3\mu_i(t)^2$
\end{itemize}

Therefore:
\begin{align}
\mathbb{E}[A_i(t)^2] &= \frac{1}{4\mu_i(t)^2} + 0 + \frac{\mu_i(t)}{4\mu_i(t)^3} + \frac{\mu_i(t)}{\mu_i(t)^2} - 0 + \frac{3\mu_i(t)^2}{4\mu_i(t)^4} \\
&= \frac{1}{4\mu_i(t)^2} + \frac{1}{4\mu_i(t)^2} + \frac{1}{\mu_i(t)} + \frac{3}{4\mu_i(t)^2} \\
&= \frac{1 + 1 + 3}{4\mu_i(t)^2} + \frac{1}{\mu_i(t)} \\
&= \frac{5}{4\mu_i(t)^2} + \frac{1}{\mu_i(t)} \\
&= \frac{5 + 4\mu_i(t)}{4\mu_i(t)^2}
\end{align}

Therefore, the Fisher Information Matrix elements are:
\begin{equation}
[\mathbf{J}]_{kl} = \sum_{t=0}^{T-1}\sum_{i=1}^{2} \frac{5 + 4\mu_i(t)}{4\mu_i(t)^2} \frac{\partial \mu_i(t)}{\partial \phi_k}\frac{\partial \mu_i(t)}{\partial \phi_l}
\end{equation}

\subsubsection{Matrix Form}

Define the Jacobian matrix:
\begin{equation}
\mathbf{G} = \frac{\partial \boldsymbol{\mu}}{\partial \boldsymbol{\phi}} \in \mathbb{R}^{2T \times 5}
\end{equation}

And the weight matrix:
\begin{equation}
\mathbf{W} = \text{diag}\left(\frac{5 + 4\mu_1(0)}{4\mu_1(0)^2}, \frac{5 + 4\mu_2(0)}{4\mu_2(0)^2}, \ldots, \frac{5 + 4\mu_2(T-1)}{4\mu_2(T-1)^2}\right)
\end{equation}

Then the Fisher Information Matrix can be written as:
\begin{equation}
\mathbf{J} = \mathbf{G}^T \mathbf{W} \mathbf{G}
\end{equation}

\subsubsection{Asymptotic Behavior}

The weight matrix elements have interesting limiting behavior:

\begin{align}
W_{ii} = \frac{5 + 4\mu_i(t)}{4\mu_i(t)^2} &= \frac{5}{4\mu_i(t)^2} + \frac{1}{\mu_i(t)} \\
&\approx \begin{cases}
\frac{5}{4\mu_i(t)^2} & \text{if } \mu_i(t) \ll 1 \text{ (small signal limit)} \\
\frac{1}{\mu_i(t)} & \text{if } \mu_i(t) \gg 1 \text{ (large signal limit)}
\end{cases}
\end{align}

\begin{itemize}
\item \textbf{Small signals}: $W_{ii} \propto 1/\mu_i(t)^2$ (similar to pure multiplicative model)
\item \textbf{Large signals}: $W_{ii} \propto 1/\mu_i(t)$ (similar to Poisson model)
\item \textbf{Transition}: The model interpolates between these regimes
\end{itemize}

\subsubsection{Cramér-Rao Lower Bound}

The Cramér-Rao lower bound for the connectivity parameter is:
\begin{equation}
\text{Var}(\hat{\theta}) \geq [\mathbf{J}^{-1}]_{11}
\end{equation}

This scaled multiplicative Gaussian model provides behavior intermediate between the pure multiplicative model (which can give unrealistically small bounds for tiny signals) and the Poisson model, while maintaining computational simplicity within the Gaussian framework.